\documentclass[11pt,a4paper]{article}
\newcommand{\gradelabel}{Grade 12}
\newcommand{\leconhdr}{Communication MQTT}
% Préambule commun aux plans de cours de robotique (Grades 7 à 12).
% Chaque plan de cours commence par \documentclass[11pt,a4paper]{article}, définit
% \gradelabel et \leconhdr (en-têtes de page), puis % Préambule commun aux plans de cours de robotique (Grades 7 à 12).
% Chaque plan de cours commence par \documentclass[11pt,a4paper]{article}, définit
% \gradelabel et \leconhdr (en-têtes de page), puis % Préambule commun aux plans de cours de robotique (Grades 7 à 12).
% Chaque plan de cours commence par \documentclass[11pt,a4paper]{article}, définit
% \gradelabel et \leconhdr (en-têtes de page), puis \input{../../preambule_robotique.tex}.
\usepackage[T1]{fontenc}
\usepackage[utf8]{inputenc}
\usepackage{lmodern}
\usepackage[french]{babel}
\usepackage[margin=2cm,headheight=14pt]{geometry}
\usepackage[most]{tcolorbox}
\usepackage{xcolor,tabularx,booktabs,array,enumitem,fancyhdr,amssymb,tikz}
\usetikzlibrary{arrows.meta,positioning}

\pagestyle{fancy}
\fancyhf{}
\lhead{\small Robotique -- \gradelabel{} -- Novembre 2026}
\rhead{\small \leconhdr}
\cfoot{\thepage}
\setlist{itemsep=1pt,topsep=2pt,leftmargin=*}
\renewcommand{\arraystretch}{1.35}
\newcolumntype{Y}{>{\raggedright\arraybackslash}X}
\newcommand{\cb}{$\square$\,}
\newcommand{\fl}{$\rightarrow$}
\newcommand{\lig}{\noindent\rule{\linewidth}{0.3pt}\\[5mm]}

\newtcolorbox{etape}[3]{enhanced,breakable,colback=#1!6,colframe=#1!70!black,
  coltitle=white,fonttitle=\bfseries,title={#2\hfill #3},boxrule=0.8pt,arc=2mm,
  before skip=8pt,after skip=8pt}
\newtcolorbox{astuce}[1][Différenciation et besoins particuliers]{breakable,
  colback=yellow!12,colframe=yellow!50!black,fonttitle=\bfseries\small,
  title={#1},boxrule=0.5pt,arc=1mm,before skip=6pt,after skip=6pt}
\newtcolorbox{code}[1][Code]{breakable,colback=black!4,colframe=black!60,
  fonttitle=\bfseries\small,title={#1},boxrule=0.5pt,arc=1mm,fontupper=\ttfamily\small,
  before skip=6pt,after skip=6pt}
% Blocs de code coloré (écrire le code sans accents : comme dans l'éditeur)
\newtcblisting{arduino}[1][Code Arduino]{listing only,breakable,colback=black!4,
  colframe=black!60,fonttitle=\bfseries\small,title={#1},boxrule=0.5pt,arc=1mm,
  before skip=6pt,after skip=6pt,
  listing options={language=C++,basicstyle=\ttfamily\small,columns=fullflexible,
  keepspaces=true,showstringspaces=false,commentstyle=\color{green!40!black}\itshape,
  keywordstyle=\color{blue!70!black}\bfseries}}
\newtcblisting{python}[1][Code Python]{listing only,breakable,colback=black!4,
  colframe=black!60,fonttitle=\bfseries\small,title={#1},boxrule=0.5pt,arc=1mm,
  before skip=6pt,after skip=6pt,
  listing options={language=Python,basicstyle=\ttfamily\small,columns=fullflexible,
  keepspaces=true,showstringspaces=false,commentstyle=\color{green!40!black}\itshape,
  keywordstyle=\color{blue!70!black}\bfseries}}
\newcommand{\entete}[2]{%
  \begin{tcolorbox}[colback=#1!10,colframe=#1!70!black,boxrule=0.8pt,arc=2mm]
  {\Large\bfseries #2}\hfill Nom : \rule{4cm}{0.3pt}\quad Date : \rule{2.2cm}{0.3pt}
  \end{tcolorbox}}

% Bandeau de titre : #1 titre, #2 semaine
\newcommand{\titreblock}[2]{%
\begin{tcolorbox}[colback=teal!10,colframe=teal!70!black,arc=2mm,boxrule=1pt]
\centering
{\LARGE\bfseries #1}\\[2mm]
{\large Robotique -- \gradelabel{} -- Plan de cours de 2 heures}\\[1mm]
Enseignant : Francis Ki \quad|\quad #2
\end{tcolorbox}}

% Tableau du minutage : 6 contenus (accroche, mini-leçon, guidée, autonome, mise en commun, exit ticket)
\newcommand{\minutage}[6]{%
\begin{tabularx}{\linewidth}{>{\bfseries}p{2.2cm}p{2.6cm}Y p{3cm}}
\toprule
Minutes & Étape & Contenu & Organisation\\
\midrule
0 -- 10 & Accroche & #1 & Classe entière\\
10 -- 30 & Mini-leçon & #2 & Classe entière\\
30 -- 60 & Pratique guidée & #3 & Collectif puis binômes\\
60 -- 95 & Pratique autonome & #4 & Individuel sur poste\\
95 -- 110 & Mise en commun & #5 & Classe entière\\
110 -- 120 & Exit ticket & #6 & Individuel sur poste\\
\bottomrule
\end{tabularx}}

% Étayage progressif : #1 aide 1, #2 aide 2, #3 aide 3
\newcommand{\aides}[3]{%
\textbf{Étayage progressif de l'aide.}
\begin{center}
\begin{tabularx}{\linewidth}{>{\bfseries}p{2.4cm}Y p{4.4cm}}
\toprule
Niveau d'aide & Contenu & Qui intervient\\
\midrule
Aide 1 & #1 & Le binôme seul\\
Aide 2 & #2 & Le binôme avec la carte d'aide\\
Aide 3 & #3 & L'enseignant\\
\bottomrule
\end{tabularx}
\end{center}}

% Les trois groupes de la pratique autonome : #1 tâche A, #2 tâche B, #3 tâche C
\newcommand{\groupes}[3]{%
\textbf{Consigne commune\,:} chaque élève ouvre le fichier de son groupe, le renomme avec son prénom et travaille seul sur son poste. L'enseignant passe d'abord auprès du groupe A pendant 10 min.

\medskip
\begin{tabularx}{\linewidth}{>{\bfseries}p{3cm}Y p{2.2cm}}
\toprule
Groupe & Tâche sur poste & Annexe\\
\midrule
A : en dessous du niveau & #1 & Annexe C\\
B : au niveau & #2 & Annexe D\\
C : au-dessus du niveau & #3 & Annexe E\\
\bottomrule
\end{tabularx}}

% Mise en commun : #1 groupe A, #2 groupe B, #3 groupe C, #4 conclusion
\newcommand{\miseencommun}[4]{%
\begin{etape}{teal}{5. Mise en commun}{95 -- 110 min}
\begin{enumerate}
 \item (2 min) Retour au calme, écrans mis en veille.
 \item (9 min) Trois passages courts de 3 min, avec valorisation :
 \begin{itemize}
  \item \textbf{Groupe A} : #1
  \item \textbf{Groupe B} : #2
  \item \textbf{Groupe C} : #3
 \end{itemize}
 \item (4 min) Conclusion : #4
\end{enumerate}
\end{etape}}

% Clôture : #1 question 1, #2 question 2
\newcommand{\cloture}[2]{%
\begin{etape}{gray}{6. Clôture : exit ticket}{110 -- 120 min}
Formulaire en ligne ou papier (annexe G) :
\begin{enumerate}
 \item #1
 \item #2
 \item Auto-évaluation : \emph{Je maîtrise / Je progresse / J'ai besoin d'aide.}
\end{enumerate}
Lire les réponses avant la séance suivante pour ajuster les groupes de besoin.
\end{etape}}

% Grille d'observation : #1 #2 #3 intitulés de colonnes
\newcommand{\grille}[3]{%
\subsection*{Grille d'observation (enseignant)}
\begin{tabularx}{\linewidth}{Y p{2.4cm} p{2.4cm} p{2.4cm} Y}
\toprule
Élève & #1 & #2 & #3 & Groupe suivant\\
\midrule
 & & & & \\ & & & & \\ & & & & \\ & & & & \\ & & & & \\
\bottomrule
\end{tabularx}}

% Signature de l'exit ticket
\newcommand{\autoeval}{Je me sens : \cb Je maîtrise \quad \cb Je progresse \quad \cb J'ai besoin d'aide}
.
\usepackage[T1]{fontenc}
\usepackage[utf8]{inputenc}
\usepackage{lmodern}
\usepackage[french]{babel}
\usepackage[margin=2cm,headheight=14pt]{geometry}
\usepackage[most]{tcolorbox}
\usepackage{xcolor,tabularx,booktabs,array,enumitem,fancyhdr,amssymb,tikz}
\usetikzlibrary{arrows.meta,positioning}

\pagestyle{fancy}
\fancyhf{}
\lhead{\small Robotique -- \gradelabel{} -- Novembre 2026}
\rhead{\small \leconhdr}
\cfoot{\thepage}
\setlist{itemsep=1pt,topsep=2pt,leftmargin=*}
\renewcommand{\arraystretch}{1.35}
\newcolumntype{Y}{>{\raggedright\arraybackslash}X}
\newcommand{\cb}{$\square$\,}
\newcommand{\fl}{$\rightarrow$}
\newcommand{\lig}{\noindent\rule{\linewidth}{0.3pt}\\[5mm]}

\newtcolorbox{etape}[3]{enhanced,breakable,colback=#1!6,colframe=#1!70!black,
  coltitle=white,fonttitle=\bfseries,title={#2\hfill #3},boxrule=0.8pt,arc=2mm,
  before skip=8pt,after skip=8pt}
\newtcolorbox{astuce}[1][Différenciation et besoins particuliers]{breakable,
  colback=yellow!12,colframe=yellow!50!black,fonttitle=\bfseries\small,
  title={#1},boxrule=0.5pt,arc=1mm,before skip=6pt,after skip=6pt}
\newtcolorbox{code}[1][Code]{breakable,colback=black!4,colframe=black!60,
  fonttitle=\bfseries\small,title={#1},boxrule=0.5pt,arc=1mm,fontupper=\ttfamily\small,
  before skip=6pt,after skip=6pt}
% Blocs de code coloré (écrire le code sans accents : comme dans l'éditeur)
\newtcblisting{arduino}[1][Code Arduino]{listing only,breakable,colback=black!4,
  colframe=black!60,fonttitle=\bfseries\small,title={#1},boxrule=0.5pt,arc=1mm,
  before skip=6pt,after skip=6pt,
  listing options={language=C++,basicstyle=\ttfamily\small,columns=fullflexible,
  keepspaces=true,showstringspaces=false,commentstyle=\color{green!40!black}\itshape,
  keywordstyle=\color{blue!70!black}\bfseries}}
\newtcblisting{python}[1][Code Python]{listing only,breakable,colback=black!4,
  colframe=black!60,fonttitle=\bfseries\small,title={#1},boxrule=0.5pt,arc=1mm,
  before skip=6pt,after skip=6pt,
  listing options={language=Python,basicstyle=\ttfamily\small,columns=fullflexible,
  keepspaces=true,showstringspaces=false,commentstyle=\color{green!40!black}\itshape,
  keywordstyle=\color{blue!70!black}\bfseries}}
\newcommand{\entete}[2]{%
  \begin{tcolorbox}[colback=#1!10,colframe=#1!70!black,boxrule=0.8pt,arc=2mm]
  {\Large\bfseries #2}\hfill Nom : \rule{4cm}{0.3pt}\quad Date : \rule{2.2cm}{0.3pt}
  \end{tcolorbox}}

% Bandeau de titre : #1 titre, #2 semaine
\newcommand{\titreblock}[2]{%
\begin{tcolorbox}[colback=teal!10,colframe=teal!70!black,arc=2mm,boxrule=1pt]
\centering
{\LARGE\bfseries #1}\\[2mm]
{\large Robotique -- \gradelabel{} -- Plan de cours de 2 heures}\\[1mm]
Enseignant : Francis Ki \quad|\quad #2
\end{tcolorbox}}

% Tableau du minutage : 6 contenus (accroche, mini-leçon, guidée, autonome, mise en commun, exit ticket)
\newcommand{\minutage}[6]{%
\begin{tabularx}{\linewidth}{>{\bfseries}p{2.2cm}p{2.6cm}Y p{3cm}}
\toprule
Minutes & Étape & Contenu & Organisation\\
\midrule
0 -- 10 & Accroche & #1 & Classe entière\\
10 -- 30 & Mini-leçon & #2 & Classe entière\\
30 -- 60 & Pratique guidée & #3 & Collectif puis binômes\\
60 -- 95 & Pratique autonome & #4 & Individuel sur poste\\
95 -- 110 & Mise en commun & #5 & Classe entière\\
110 -- 120 & Exit ticket & #6 & Individuel sur poste\\
\bottomrule
\end{tabularx}}

% Étayage progressif : #1 aide 1, #2 aide 2, #3 aide 3
\newcommand{\aides}[3]{%
\textbf{Étayage progressif de l'aide.}
\begin{center}
\begin{tabularx}{\linewidth}{>{\bfseries}p{2.4cm}Y p{4.4cm}}
\toprule
Niveau d'aide & Contenu & Qui intervient\\
\midrule
Aide 1 & #1 & Le binôme seul\\
Aide 2 & #2 & Le binôme avec la carte d'aide\\
Aide 3 & #3 & L'enseignant\\
\bottomrule
\end{tabularx}
\end{center}}

% Les trois groupes de la pratique autonome : #1 tâche A, #2 tâche B, #3 tâche C
\newcommand{\groupes}[3]{%
\textbf{Consigne commune\,:} chaque élève ouvre le fichier de son groupe, le renomme avec son prénom et travaille seul sur son poste. L'enseignant passe d'abord auprès du groupe A pendant 10 min.

\medskip
\begin{tabularx}{\linewidth}{>{\bfseries}p{3cm}Y p{2.2cm}}
\toprule
Groupe & Tâche sur poste & Annexe\\
\midrule
A : en dessous du niveau & #1 & Annexe C\\
B : au niveau & #2 & Annexe D\\
C : au-dessus du niveau & #3 & Annexe E\\
\bottomrule
\end{tabularx}}

% Mise en commun : #1 groupe A, #2 groupe B, #3 groupe C, #4 conclusion
\newcommand{\miseencommun}[4]{%
\begin{etape}{teal}{5. Mise en commun}{95 -- 110 min}
\begin{enumerate}
 \item (2 min) Retour au calme, écrans mis en veille.
 \item (9 min) Trois passages courts de 3 min, avec valorisation :
 \begin{itemize}
  \item \textbf{Groupe A} : #1
  \item \textbf{Groupe B} : #2
  \item \textbf{Groupe C} : #3
 \end{itemize}
 \item (4 min) Conclusion : #4
\end{enumerate}
\end{etape}}

% Clôture : #1 question 1, #2 question 2
\newcommand{\cloture}[2]{%
\begin{etape}{gray}{6. Clôture : exit ticket}{110 -- 120 min}
Formulaire en ligne ou papier (annexe G) :
\begin{enumerate}
 \item #1
 \item #2
 \item Auto-évaluation : \emph{Je maîtrise / Je progresse / J'ai besoin d'aide.}
\end{enumerate}
Lire les réponses avant la séance suivante pour ajuster les groupes de besoin.
\end{etape}}

% Grille d'observation : #1 #2 #3 intitulés de colonnes
\newcommand{\grille}[3]{%
\subsection*{Grille d'observation (enseignant)}
\begin{tabularx}{\linewidth}{Y p{2.4cm} p{2.4cm} p{2.4cm} Y}
\toprule
Élève & #1 & #2 & #3 & Groupe suivant\\
\midrule
 & & & & \\ & & & & \\ & & & & \\ & & & & \\ & & & & \\
\bottomrule
\end{tabularx}}

% Signature de l'exit ticket
\newcommand{\autoeval}{Je me sens : \cb Je maîtrise \quad \cb Je progresse \quad \cb J'ai besoin d'aide}
.
\usepackage[T1]{fontenc}
\usepackage[utf8]{inputenc}
\usepackage{lmodern}
\usepackage[french]{babel}
\usepackage[margin=2cm,headheight=14pt]{geometry}
\usepackage[most]{tcolorbox}
\usepackage{xcolor,tabularx,booktabs,array,enumitem,fancyhdr,amssymb,tikz}
\usetikzlibrary{arrows.meta,positioning}

\pagestyle{fancy}
\fancyhf{}
\lhead{\small Robotique -- \gradelabel{} -- Novembre 2026}
\rhead{\small \leconhdr}
\cfoot{\thepage}
\setlist{itemsep=1pt,topsep=2pt,leftmargin=*}
\renewcommand{\arraystretch}{1.35}
\newcolumntype{Y}{>{\raggedright\arraybackslash}X}
\newcommand{\cb}{$\square$\,}
\newcommand{\fl}{$\rightarrow$}
\newcommand{\lig}{\noindent\rule{\linewidth}{0.3pt}\\[5mm]}

\newtcolorbox{etape}[3]{enhanced,breakable,colback=#1!6,colframe=#1!70!black,
  coltitle=white,fonttitle=\bfseries,title={#2\hfill #3},boxrule=0.8pt,arc=2mm,
  before skip=8pt,after skip=8pt}
\newtcolorbox{astuce}[1][Différenciation et besoins particuliers]{breakable,
  colback=yellow!12,colframe=yellow!50!black,fonttitle=\bfseries\small,
  title={#1},boxrule=0.5pt,arc=1mm,before skip=6pt,after skip=6pt}
\newtcolorbox{code}[1][Code]{breakable,colback=black!4,colframe=black!60,
  fonttitle=\bfseries\small,title={#1},boxrule=0.5pt,arc=1mm,fontupper=\ttfamily\small,
  before skip=6pt,after skip=6pt}
% Blocs de code coloré (écrire le code sans accents : comme dans l'éditeur)
\newtcblisting{arduino}[1][Code Arduino]{listing only,breakable,colback=black!4,
  colframe=black!60,fonttitle=\bfseries\small,title={#1},boxrule=0.5pt,arc=1mm,
  before skip=6pt,after skip=6pt,
  listing options={language=C++,basicstyle=\ttfamily\small,columns=fullflexible,
  keepspaces=true,showstringspaces=false,commentstyle=\color{green!40!black}\itshape,
  keywordstyle=\color{blue!70!black}\bfseries}}
\newtcblisting{python}[1][Code Python]{listing only,breakable,colback=black!4,
  colframe=black!60,fonttitle=\bfseries\small,title={#1},boxrule=0.5pt,arc=1mm,
  before skip=6pt,after skip=6pt,
  listing options={language=Python,basicstyle=\ttfamily\small,columns=fullflexible,
  keepspaces=true,showstringspaces=false,commentstyle=\color{green!40!black}\itshape,
  keywordstyle=\color{blue!70!black}\bfseries}}
\newcommand{\entete}[2]{%
  \begin{tcolorbox}[colback=#1!10,colframe=#1!70!black,boxrule=0.8pt,arc=2mm]
  {\Large\bfseries #2}\hfill Nom : \rule{4cm}{0.3pt}\quad Date : \rule{2.2cm}{0.3pt}
  \end{tcolorbox}}

% Bandeau de titre : #1 titre, #2 semaine
\newcommand{\titreblock}[2]{%
\begin{tcolorbox}[colback=teal!10,colframe=teal!70!black,arc=2mm,boxrule=1pt]
\centering
{\LARGE\bfseries #1}\\[2mm]
{\large Robotique -- \gradelabel{} -- Plan de cours de 2 heures}\\[1mm]
Enseignant : Francis Ki \quad|\quad #2
\end{tcolorbox}}

% Tableau du minutage : 6 contenus (accroche, mini-leçon, guidée, autonome, mise en commun, exit ticket)
\newcommand{\minutage}[6]{%
\begin{tabularx}{\linewidth}{>{\bfseries}p{2.2cm}p{2.6cm}Y p{3cm}}
\toprule
Minutes & Étape & Contenu & Organisation\\
\midrule
0 -- 10 & Accroche & #1 & Classe entière\\
10 -- 30 & Mini-leçon & #2 & Classe entière\\
30 -- 60 & Pratique guidée & #3 & Collectif puis binômes\\
60 -- 95 & Pratique autonome & #4 & Individuel sur poste\\
95 -- 110 & Mise en commun & #5 & Classe entière\\
110 -- 120 & Exit ticket & #6 & Individuel sur poste\\
\bottomrule
\end{tabularx}}

% Étayage progressif : #1 aide 1, #2 aide 2, #3 aide 3
\newcommand{\aides}[3]{%
\textbf{Étayage progressif de l'aide.}
\begin{center}
\begin{tabularx}{\linewidth}{>{\bfseries}p{2.4cm}Y p{4.4cm}}
\toprule
Niveau d'aide & Contenu & Qui intervient\\
\midrule
Aide 1 & #1 & Le binôme seul\\
Aide 2 & #2 & Le binôme avec la carte d'aide\\
Aide 3 & #3 & L'enseignant\\
\bottomrule
\end{tabularx}
\end{center}}

% Les trois groupes de la pratique autonome : #1 tâche A, #2 tâche B, #3 tâche C
\newcommand{\groupes}[3]{%
\textbf{Consigne commune\,:} chaque élève ouvre le fichier de son groupe, le renomme avec son prénom et travaille seul sur son poste. L'enseignant passe d'abord auprès du groupe A pendant 10 min.

\medskip
\begin{tabularx}{\linewidth}{>{\bfseries}p{3cm}Y p{2.2cm}}
\toprule
Groupe & Tâche sur poste & Annexe\\
\midrule
A : en dessous du niveau & #1 & Annexe C\\
B : au niveau & #2 & Annexe D\\
C : au-dessus du niveau & #3 & Annexe E\\
\bottomrule
\end{tabularx}}

% Mise en commun : #1 groupe A, #2 groupe B, #3 groupe C, #4 conclusion
\newcommand{\miseencommun}[4]{%
\begin{etape}{teal}{5. Mise en commun}{95 -- 110 min}
\begin{enumerate}
 \item (2 min) Retour au calme, écrans mis en veille.
 \item (9 min) Trois passages courts de 3 min, avec valorisation :
 \begin{itemize}
  \item \textbf{Groupe A} : #1
  \item \textbf{Groupe B} : #2
  \item \textbf{Groupe C} : #3
 \end{itemize}
 \item (4 min) Conclusion : #4
\end{enumerate}
\end{etape}}

% Clôture : #1 question 1, #2 question 2
\newcommand{\cloture}[2]{%
\begin{etape}{gray}{6. Clôture : exit ticket}{110 -- 120 min}
Formulaire en ligne ou papier (annexe G) :
\begin{enumerate}
 \item #1
 \item #2
 \item Auto-évaluation : \emph{Je maîtrise / Je progresse / J'ai besoin d'aide.}
\end{enumerate}
Lire les réponses avant la séance suivante pour ajuster les groupes de besoin.
\end{etape}}

% Grille d'observation : #1 #2 #3 intitulés de colonnes
\newcommand{\grille}[3]{%
\subsection*{Grille d'observation (enseignant)}
\begin{tabularx}{\linewidth}{Y p{2.4cm} p{2.4cm} p{2.4cm} Y}
\toprule
Élève & #1 & #2 & #3 & Groupe suivant\\
\midrule
 & & & & \\ & & & & \\ & & & & \\ & & & & \\ & & & & \\
\bottomrule
\end{tabularx}}

% Signature de l'exit ticket
\newcommand{\autoeval}{Je me sens : \cb Je maîtrise \quad \cb Je progresse \quad \cb J'ai besoin d'aide}


\begin{document}
\titreblock{Communication MQTT}{Semaine du 23 au 27 novembre 2026}

%==================================================================
\section{Informations générales}
\begin{tabularx}{\linewidth}{>{\bfseries}p{3.6cm}Y}
\toprule
Titre & Communication MQTT : publier et recevoir des données\\
Classe / effectif & Grade 12, 11 à 20 élèves, niveaux hétérogènes, élèves à besoins particuliers\\
Durée & 2 heures (120 min)\\
Lieu & Salle informatique, un poste par élève\\
Matériel et logiciels &
\begin{itemize}
 \item Un poste par élève avec navigateur web.
 \item \textbf{Wokwi} (wokwi.com, gratuit) : ESP32 simulé avec Wi-Fi \texttt{Wokwi-GUEST}, bibliothèque \texttt{PubSubClient}, un capteur de température et une LED.
 \item Un client MQTT web gratuit (par exemple le client WebSocket de HiveMQ) et un courtier public d'essai (par exemple \texttt{broker.hivemq.com}). \textbf{À tester sur un poste élève avant la séance.}
 \item Facultatif : une carte ESP32 réelle par binôme.
 \item Vidéoprojecteur ou tableau interactif.
\end{itemize}\\
Pré-requis &
Séances précédentes : architecture d'un objet connecté, protocoles (HTTP, MQTT), connexion Wi-Fi de l'ESP32 avec adresse IP et gestion d'échec.\\
\bottomrule
\end{tabularx}

\begin{astuce}[Préparation avant le cours (environ 55 min)]
\begin{enumerate}
 \item Créer le projet Wokwi (ESP32, capteur de température, LED) avec la bibliothèque PubSubClient, et vérifier la connexion au courtier public depuis un poste élève. Si le pare-feu de l'école bloque le port MQTT (1883), prévoir le plan B.
 \item Définir un préfixe de sujets unique à la classe, par exemple \texttt{ibsa/g12/}, suivi du prénom de l'élève.
 \item Expliquer et faire respecter la règle : \textbf{aucune donnée personnelle} sur un courtier public. Tout le monde peut lire les messages.
 \item Préparer les fichiers des annexes C à E et le formulaire de sortie.
\end{enumerate}
\end{astuce}

%==================================================================
\section{Objectifs d'apprentissage}
\begin{tcolorbox}[colback=gray!8,colframe=gray!60!black,arc=1mm,boxrule=0.5pt]
\textbf{Objectifs du programme}
\begin{itemize}
 \item Comprendre le protocole MQTT.
 \item Publier et recevoir des données.
\end{itemize}
\textbf{À la fin de la séance, l'élève est capable\,:}
\begin{itemize}
 \item d'expliquer le rôle du courtier (broker), de l'éditeur (publisher), de l'abonné (subscriber) et du sujet (topic)\,;
 \item de publier une mesure sur un sujet et de s'abonner à un sujet pour recevoir une commande\,;
 \item d'organiser des sujets de façon claire et de justifier le niveau de sécurité d'un échange.
\end{itemize}
\textbf{Compétences numériques\,:} publier et s'abonner, lire un journal d'échanges, protéger ses données.
\end{tcolorbox}

%==================================================================
\section{Déroulé minuté (120 min)}
\minutage{Le journal et ses abonnés}{Courtier, sujet, publication, abonnement}{Publication et abonnement avec le client web en collectif, puis avec l'ESP32 simulé en binômes}{Trois niveaux de besoin}{Valoriser chaque groupe}{Trois questions sur poste}

%------------------------------------------------------------------
\begin{etape}{orange}{1. Accroche : le journal et ses abonnés}{0 -- 10 min}
\begin{enumerate}
 \item (5 min) Mise en situation : un élève «~imprimeur~» (courtier) reçoit des articles de plusieurs «~journalistes~» sur des rubriques (sport, météo) et les distribue aux «~abonnés~» de chaque rubrique. Les journalistes ne connaissent pas les abonnés.
 \item (5 min) Question-problème : \textbf{«~Comment un capteur peut-il envoyer sa mesure à plusieurs applications sans les connaître ?~»}
\end{enumerate}
\end{etape}

%------------------------------------------------------------------
\begin{etape}{blue}{2. Mini-leçon : le principe de MQTT}{10 -- 30 min}
\textbf{Support\,:} annexe A. Quatre temps de 5 min.

\textbf{Temps 1 : trois acteurs.}
\begin{center}
\begin{tikzpicture}[box/.style={draw,rounded corners,thick,minimum height=1.1cm,text width=2.6cm,align=center,font=\small},>={Stealth[length=3mm]},node distance=1cm]
\node[box,fill=orange!20] (p) {Éditeur\\ ESP32 + capteur};
\node[box,fill=blue!15,right=of p] (b) {Courtier\\ (broker)};
\node[box,fill=green!15,right=of b] (s) {Abonné\\ application};
\draw[->,thick] (p)--node[above,font=\scriptsize]{publie}(b);
\draw[->,thick] (b)--node[above,font=\scriptsize]{transmet}(s);
\end{tikzpicture}
\end{center}

\textbf{Temps 2 : les sujets (topics).} Hiérarchiques, séparés par des barres obliques : \texttt{ibsa/g12/awa/temp}. Deux jokers : \texttt{+} (un niveau) et \texttt{\#} (tous les niveaux suivants). Exemple : s'abonner à \texttt{ibsa/g12/+/temp} reçoit la température de tous les élèves.

\textbf{Temps 3 : le programme (PubSubClient).}
\begin{arduino}[Publier et s'abonner (extrait)]
#include <WiFi.h>
#include <PubSubClient.h>

WiFiClient reseau;
PubSubClient client(reseau);

void recevoir(char* sujet, byte* donnees, unsigned int taille) {
  String message = "";
  for (unsigned int i = 0; i < taille; i++) {
    message += (char)donnees[i];
  }
  Serial.println(message);
  if (message == "on") digitalWrite(5, HIGH);
  if (message == "off") digitalWrite(5, LOW);
}

void setup() {
  // connexion Wi-Fi comme a la seance precedente
  client.setServer("broker.hivemq.com", 1883);
  client.setCallback(recevoir);
}

void loop() {
  if (!client.connected()) {
    client.connect("ibsa-g12-awa");
    client.subscribe("ibsa/g12/awa/led");
  }
  client.loop();
  client.publish("ibsa/g12/awa/temp", "24");
  delay(5000);
}
\end{arduino}

\textbf{Temps 4 : sécurité.} Un courtier public n'est ni chiffré ni protégé : pas de mots de passe, pas de données personnelles. En production : authentification, chiffrement TLS (port 8883) et contrôle d'accès par sujet.

\begin{astuce}
Colorer éditeur, courtier et abonné dans les schémas et dans le code. Lire le code en disant «~publie sur le sujet\ldots~» et «~s'abonne au sujet\ldots~». Fournir le lexique anglais (publish, subscribe, topic, broker).
\end{astuce}
\end{etape}

%------------------------------------------------------------------
\begin{etape}{green}{3. Pratique guidée}{30 -- 60 min}
\textbf{Fiche\,:} annexe B. Programmation en binômes (pilote et navigateur, rôles échangés à mi-temps).

\textbf{3a. Client web en collectif (12 min).} Dans le client MQTT web, tous se connectent au courtier, s'abonnent à \texttt{ibsa/g12/+/temp}, puis chaque élève publie sa valeur sur \texttt{ibsa/g12/prenom/temp}. Lire les messages reçus.

\textbf{3b. ESP32 simulé en binômes (18 min).} Ouvrir le projet Wokwi, remplacer le prénom dans les sujets, lancer la simulation, observer la publication de la température dans le client web. Publier «~on~» puis «~off~» sur le sujet de la LED depuis le client web et vérifier la réaction. Alternance des rôles.

\aides{Le tableau des sujets de la classe et le programme annoté}{Questions guides : \emph{Qui est l'éditeur ? Qui est l'abonné ? Sur quel sujet publie-t-on ?}}{L'enseignant vérifie le sujet avec le binôme et relance la connexion}

\begin{astuce}
Contrôler que chaque sujet contient bien le prénom de l'élève (éviter les collisions entre binômes). Si le courtier est inaccessible, passer au plan B (jeu de rôles avec le client web seul ou avec cartes sujets).
\end{astuce}
\end{etape}

%------------------------------------------------------------------
\begin{etape}{violet}{4. Pratique autonome différenciée}{60 -- 95 min}
\groupes{Programme à compléter (instructions manquantes : \texttt{setServer}, \texttt{connect}, \texttt{subscribe}, \texttt{publish}, banque de mots). Relier les rôles (éditeur, abonné, courtier) à leurs éléments dans un schéma.}{Programmer seul une station qui publie la température toutes les 5 s sur \texttt{ibsa/g12/prenom/temp} et qui allume la LED à la réception de «~on~» sur \texttt{ibsa/g12/prenom/led}. Fournir un tableau de tests.}{Concevoir une organisation de sujets pour une station complète (température, humidité, état, commande), utiliser les jokers \texttt{+} et \texttt{\#} pour un tableau de bord, et rédiger une analyse de sécurité du système (courtier public, authentification, TLS).}

\medskip
\textbf{Plan B si le courtier est bloqué\,:} faire jouer le courtier à un élève et tenir un journal papier des publications et abonnements.
\end{etape}

%------------------------------------------------------------------
\miseencommun{un élève montre la valeur reçue dans le client web et explique qui publie. Valoriser l'exactitude.}{un binôme montre la commande de la LED depuis le client web. Valoriser le lien entre le sujet et l'action.}{un élève présente son organisation de sujets et son analyse de sécurité. Valoriser la réflexion sur les risques.}{\textit{MQTT sépare ceux qui publient de ceux qui reçoivent grâce à un courtier et à des sujets ; il est léger, mais il faut le sécuriser.}}

\cloture{Quel est le rôle du courtier (broker) dans MQTT ?}{Pourquoi ne faut-il pas envoyer de données personnelles sur un courtier public ?}

%==================================================================
% ANNEXES
%==================================================================
\newpage
\section*{Annexe A : Support résumé de la mini-leçon}
\entete{blue}{Communication MQTT}

\begin{tcolorbox}[colback=yellow!10,colframe=yellow!50!black]
\textbf{Éditeur} publie sur un \textbf{sujet}. \textbf{Courtier} transmet. \textbf{Abonné} reçoit les messages des sujets auxquels il est abonné.
\end{tcolorbox}

\begin{center}
\begin{tikzpicture}[box/.style={draw,rounded corners,thick,minimum height=1.1cm,text width=2.6cm,align=center,font=\small},>={Stealth[length=3mm]},node distance=1cm]
\node[box,fill=orange!20] (p) {Éditeur};
\node[box,fill=blue!15,right=of p] (b) {Courtier};
\node[box,fill=green!15,right=of b] (s) {Abonné};
\draw[->,thick] (p)--(b); \draw[->,thick] (b)--(s);
\end{tikzpicture}
\end{center}

\begin{tabularx}{\linewidth}{>{\ttfamily}p{4.4cm}Y}
\toprule
Instruction & Rôle\\
\midrule
client.setServer(hote, 1883) & Adresse du courtier\\
client.connect("identifiant") & Se connecte (identifiant unique)\\
client.subscribe("sujet") & S'abonne\\
client.publish("sujet", "valeur") & Publie\\
client.loop() & Maintient la connexion, reçoit\\
\bottomrule
\end{tabularx}

\medskip
\textbf{Jokers :} \texttt{+} (un niveau), \texttt{\#} (tous les niveaux suivants).\\
\textbf{Règle de la classe :} aucun mot de passe, aucune donnée personnelle sur un courtier public.

%==================================================================
\newpage
\section*{Annexe B : Fiche de pratique guidée}
\entete{green}{Pratique guidée : publier et recevoir}

\subsection*{3a. Client web}
Mon sujet : \dotfill\\
Valeur publiée : \dotfill\\
Messages reçus (au moins trois) : \dotfill

\subsection*{3b. ESP32 simulé}
\begin{tabularx}{\linewidth}{>{\bfseries}p{3.6cm}Y Y}
\toprule
Action & Sujet utilisé & Résultat observé\\
\midrule
L'ESP32 publie la température & & \\
Je publie «~on~» & & \\
Je publie «~off~» & & \\
\bottomrule
\end{tabularx}

\begin{tcolorbox}[colback=yellow!10,colframe=yellow!50!black,title={Cartes d'aide à utiliser dans l'ordre},fonttitle=\bfseries\small]
\textbf{Aide 1 :} regarde le tableau des sujets et le programme annoté.\\
\textbf{Aide 2 :} Qui est l'éditeur ? Qui est l'abonné ? Sur quel sujet publie-t-on ?\\
\textbf{Aide 3 :} demande à l'enseignant de vérifier ton sujet avec toi.
\end{tcolorbox}

%==================================================================
\newpage
\section*{Annexe C : Niveau A (en dessous du niveau)}
\entete{orange}{Niveau A : je complète le programme MQTT}

\textbf{Banque :} setServer, connect, subscribe, publish.

\begin{arduino}[Programme à compléter]
client._____("broker.hivemq.com", 1883);
client._____("ibsa-g12-awa");
client._____("ibsa/g12/awa/led");
client._____("ibsa/g12/awa/temp", "24");
\end{arduino}

\subsection*{Activité : relie}
\begin{tabularx}{\linewidth}{Y Y}
\toprule
\textbf{Rôle} & \textbf{Élément}\\
\midrule
Éditeur & \\
Courtier & \\
Abonné & \\
\bottomrule
\end{tabularx}
Éléments : ESP32 avec capteur, serveur broker.hivemq.com, application qui affiche la température.

\subsection*{Vrai ou faux ?}
1. Un abonné doit connaître l'éditeur. \cb Vrai \cb Faux\\
2. Un sujet ressemble à un chemin avec des barres obliques. \cb Vrai \cb Faux

%==================================================================
\newpage
\section*{Annexe D : Niveau B (au niveau)}
\entete{blue}{Niveau B : ma station MQTT}

\textbf{Cahier des charges :} publier la température toutes les 5 s sur \texttt{ibsa/g12/prenom/temp} ; allumer la LED à la réception de «~on~» et l'éteindre à la réception de «~off~» sur \texttt{ibsa/g12/prenom/led}.

\begin{tabularx}{\linewidth}{>{\bfseries}p{4cm}Y Y Y}
\toprule
Test & Résultat attendu & Obtenu & Conforme ?\\
\midrule
Publication de la température & & & \\
Commande «~on~» & & & \\
Commande «~off~» & & & \\
Message inconnu («~xyz~») & & & \\
\bottomrule
\end{tabularx}

\subsection*{Question}
Que se passe-t-il si deux élèves utilisent le même identifiant de client ? \dotfill

%==================================================================
\newpage
\section*{Annexe E : Niveau C (au-dessus du niveau)}
\entete{violet}{Niveau C : sujets et sécurité}

\subsection*{1. Organisation des sujets d'une station complète}
\begin{tabularx}{\linewidth}{>{\bfseries}p{3.6cm}Y Y}
\toprule
Information & Sujet & Qui publie ?\\
\midrule
Température & & \\ Humidité & & \\ État de la station & & \\ Commande d'arrosage & & \\
\bottomrule
\end{tabularx}

\subsection*{2. Tableau de bord}
Quel abonnement permet de recevoir toutes les mesures de la station ? \dotfill\\
Quel abonnement permet de recevoir la température de toutes les stations ? \dotfill

\subsection*{3. Analyse de sécurité}
\begin{tabularx}{\linewidth}{>{\bfseries}p{3.6cm}Y Y}
\toprule
Risque & Exemple & Parade\\
\midrule
Lecture des messages & & \\ Commande par un inconnu & & \\ Usurpation d'identité & & \\
\bottomrule
\end{tabularx}

\textbf{Défi bonus :} explique la différence entre les niveaux de qualité de service (0, 1, 2) de MQTT.

%==================================================================
\newpage
\section*{Annexe F : Corrections et attendus (enseignant)}

\subsection*{Annexe B}
Publication : \texttt{ibsa/g12/prenom/temp}. Commande : \texttt{ibsa/g12/prenom/led}. La LED s'allume avec «~on~» et s'éteint avec «~off~».

\subsection*{Annexe C}
\texttt{setServer}, \texttt{connect}, \texttt{subscribe}, \texttt{publish}. Rôles : ESP32 = éditeur ; broker = courtier ; application = abonné. Faux, Vrai.

\subsection*{Annexe D}
Message inconnu : la LED ne change pas. Deux clients avec le même identifiant se déconnectent mutuellement (le courtier garde le dernier).

\subsection*{Annexe E}
Sujets : \texttt{ibsa/g12/prenom/temp}, \texttt{.../hum}, \texttt{.../etat}, \texttt{.../arrosage}. Abonnements : \texttt{ibsa/g12/prenom/\#} ; \texttt{ibsa/g12/+/temp}. Risques : lecture (chiffrer avec TLS), commande par un inconnu (authentification et contrôle d'accès), usurpation (identifiant unique et mot de passe). QoS : 0 au plus une fois, 1 au moins une fois, 2 exactement une fois.\\
\textbf{Critères de réussite :} sujets cohérents, jokers corrects, trois risques et trois parades.

%==================================================================
\newpage
\section*{Annexe G : Exit ticket et grille d'observation}
\entete{gray}{Exit ticket : MQTT}

\begin{enumerate}
 \item Quel est le rôle du courtier (broker) dans MQTT ?\\ \lig
 \item Pourquoi ne faut-il pas envoyer de données personnelles sur un courtier public ?\\ \lig
 \item \autoeval
\end{enumerate}
\textbf{Attendus :} 1. Recevoir les messages des éditeurs et les transmettre aux abonnés du sujet. 2. Tout le monde peut lire les messages (non chiffrés, non protégés).

\bigskip
\grille{Publication}{Abonnement}{Sécurité comprise}

\end{document}
